\scanpage{87}
\mspage{539}

\noindent\qquad 3.8. Quasi-homéomorphismes et images réciproques de faisceaux\dotfill

\noindent\qquad 3.9. Faisceaux d'espaces pseudo-discrets\dotfill

\noindent\quad \S\ 4. Espaces annelés\dotfill

\noindent\qquad 4.1. Espaces annelés , $A$-Modules , $A$-Algèbres\dotfill

\noindent\qquad 4.2. Image directe d'un $A$-Module\dotfill

\noindent\qquad 4.3. Image réciproque d'un $B$-Module\dotfill

\noindent\qquad 4.4. Relations entre images directes et images réciproques\dotfill

\noindent\qquad 4.5.Immersions ouvertes et morphismes représentables par des immersions ouvertes \dotfill

\noindent\quad \S\ 5. Faisceaux quasi-cohérents et faisceaux cohérents\dotfill

\noindent\qquad 5.1. Faisceaux quasi-cohérents\dotfill

\noindent\qquad 5.2. Faisceaux de type fini\dotfill

\noindent\qquad 5.3. Faisceaux cohérents\dotfill

\noindent\qquad 5.4. Faisceaux localement libres\dotfill

\noindent\qquad 5.5. Modules localement libres sur un espace localement annelé\dotfill

\noindent\qquad 5.6. Groupe de Picard\dotfill

\noindent\qquad 5.7. Morphismes plats d'espaces annelés\dotfill

\noindent\quad \S\ 6. Compléments d'algèbre commutative .

\noindent\qquad 6.1. Limites inductives d'anneaux\dotfill

\noindent\qquad 6.2. Identification du module $M_f$ à une limite \deleted{\unsure{project|projecti}} inductive\dotfill

\noindent\qquad 6.3. Propriétés de finitude\dotfill

\noindent\qquad 6.4. Critères pour qu'un anneau soit noethérien\dotfill

\noindent\qquad 6.5. Anneaux normaux et ferm\unsure{\deleted{t}e\deleted{e}t|\deleted{e}t\deleted{t}e}ure intégrale\dotfill

\noindent\qquad 6.6. Compléments sur la platitude\dotfill

\noindent\qquad 6.7. Compléments sur\deleted{l} les modules projectifs\dotfill

\noindent\qquad 6.8. Existence d'extensions plates d'anneaux locaux\dotfill

\noindent\quad \deleted{\S\ \unsure{7.}}

\noindent\quad \S\ 7. Compléments d'algèbre topologique\dotfill

\noindent\qquad 7.1. Anneaux admissibles\dotfill

\noindent\qquad 7.2. Anneaux adiques et limites projectives\dotfill

\noindent\qquad 7.3. Anneaux préadiques noethériens\dotfill

\noindent\qquad 7.4. Modules quasi-finis sur les anneaux locaux\dotfill

\noindent\qquad \deleted{7\unsure{.}} 7.5. Anneaux de séries formelles restreintes\dotfill

\noindent\qquad 7.6. Anneaux complets de fractions\dotfill

% archive sheet number: 86
